% Shared Beamer preamble for AI Frameworks
% Audience: college students ages 16--18 (intuition over math)
% Include with: % Shared Beamer preamble for AI Frameworks
% Audience: college students ages 16--18 (intuition over math)
% Include with: % Shared Beamer preamble for AI Frameworks
% Audience: college students ages 16--18 (intuition over math)
% Include with: % Shared Beamer preamble for AI Frameworks
% Audience: college students ages 16--18 (intuition over math)
% Include with: \input{preamble}

\usetheme{Madrid}
\usecolortheme{default}

% Clean palette: deep slate + teal accent (distinct from Java navy decks)
\definecolor{LectureNavy}{HTML}{1A365D}
\definecolor{LectureAccent}{HTML}{319795}
\definecolor{LectureGray}{HTML}{4A5568}
\definecolor{CodeBg}{HTML}{F7FAFC}
\definecolor{AlertBg}{HTML}{FFF5F5}
\definecolor{TryBg}{HTML}{F0FFF4}
\definecolor{QuizBg}{HTML}{EBF8FF}
\definecolor{AnswerKeyBg}{HTML}{FFFFF0}
\definecolor{AnswerGold}{HTML}{B7791F}
\definecolor{KeywordBlue}{HTML}{2B6CB0}
\definecolor{StringGreen}{HTML}{276749}
\definecolor{CommentGray}{HTML}{718096}

\setbeamercolor{structure}{fg=LectureNavy}
\setbeamercolor{palette primary}{bg=LectureNavy,fg=white}
\setbeamercolor{palette secondary}{bg=LectureAccent,fg=white}
\setbeamercolor{palette tertiary}{bg=LectureGray,fg=white}
% Madrid title box is dark --- title/subtitle must be light for contrast
\setbeamercolor{title}{fg=white,bg=LectureNavy}
\setbeamercolor{subtitle}{fg=white,bg=LectureNavy}
\setbeamercolor{author}{fg=LectureNavy}
\setbeamercolor{date}{fg=LectureGray}
\setbeamercolor{institute}{fg=LectureGray}
\setbeamercolor{frametitle}{bg=LectureNavy,fg=white}
\setbeamercolor{block title}{bg=LectureNavy,fg=white}
\setbeamercolor{block body}{bg=CodeBg,fg=black}
\setbeamercolor{block title alerted}{bg=LectureAccent,fg=white}
\setbeamercolor{block body alerted}{bg=TryBg,fg=black}
\setbeamercolor{block title example}{bg=LectureGray,fg=white}
\setbeamercolor{block body example}{bg=AlertBg,fg=black}

\setbeamertemplate{navigation symbols}{}
\setbeamertemplate{itemize items}[circle]
\setbeamertemplate{enumerate items}[default]

\usepackage[T1]{fontenc}
\usepackage[utf8]{inputenc}
\usepackage{lmodern}
\usepackage{booktabs}
\usepackage{tabularx}
\usepackage{graphicx}
\usepackage{hyperref}
\usepackage{listings}
\usepackage{xcolor}
\usepackage{tikz}
\usetikzlibrary{positioning,arrows.meta,shapes.geometric}

% listings: Python without shell-escape
\lstdefinestyle{pythonstyle}{
    language=Python,
    basicstyle=\ttfamily\small,
    keywordstyle=\color{KeywordBlue}\bfseries,
    commentstyle=\color{CommentGray}\itshape,
    stringstyle=\color{StringGreen},
    numbers=left,
    numberstyle=\tiny\color{CommentGray},
    numbersep=6pt,
    frame=single,
    rulecolor=\color{LectureGray!40},
    backgroundcolor=\color{CodeBg},
    breaklines=true,
    showstringspaces=false,
    tabsize=4,
    captionpos=b,
}
\lstset{style=pythonstyle}

\newcommand{\pycode}[1]{{\ttfamily\small\detokenize{#1}}}
\newcommand{\trythis}[1]{%
    \begin{alertblock}{Try this}
    #1
    \end{alertblock}%
}
\newcommand{\commonmistake}[1]{%
    \begin{exampleblock}{Common mistake}
    #1
    \end{exampleblock}%
}
\newcommand{\remember}[1]{%
    \begin{block}{Remember}
    #1
    \end{block}%
}

% Formative in-class checks: question slide, then a dedicated Answer slide
\newcommand{\quizpause}{%
    \par\vspace{0.5em}
    {\small\textit{Pause --- answer (clicker / raise hand / neighbour) before the next slide.}}%
}
\newcommand{\quizbox}[1]{%
    \setbeamercolor{block title}{bg=KeywordBlue,fg=white}%
    \setbeamercolor{block body}{bg=QuizBg,fg=black}%
    \begin{block}{Quiz}
    #1
    \end{block}%
    \setbeamercolor{block title}{bg=LectureNavy,fg=white}%
    \setbeamercolor{block body}{bg=CodeBg,fg=black}%
}
\newcommand{\answerbox}[1]{%
    \setbeamercolor{block title}{bg=AnswerGold,fg=white}%
    \setbeamercolor{block body}{bg=AnswerKeyBg,fg=black}%
    \begin{block}{Answer}
    #1
    \end{block}%
    \setbeamercolor{block title}{bg=LectureNavy,fg=white}%
    \setbeamercolor{block body}{bg=CodeBg,fg=black}%
}

% Empty author/date placeholders for the lecturer to fill
\author{}
\date{}


\usetheme{Madrid}
\usecolortheme{default}

% Clean palette: deep slate + teal accent (distinct from Java navy decks)
\definecolor{LectureNavy}{HTML}{1A365D}
\definecolor{LectureAccent}{HTML}{319795}
\definecolor{LectureGray}{HTML}{4A5568}
\definecolor{CodeBg}{HTML}{F7FAFC}
\definecolor{AlertBg}{HTML}{FFF5F5}
\definecolor{TryBg}{HTML}{F0FFF4}
\definecolor{QuizBg}{HTML}{EBF8FF}
\definecolor{AnswerKeyBg}{HTML}{FFFFF0}
\definecolor{AnswerGold}{HTML}{B7791F}
\definecolor{KeywordBlue}{HTML}{2B6CB0}
\definecolor{StringGreen}{HTML}{276749}
\definecolor{CommentGray}{HTML}{718096}

\setbeamercolor{structure}{fg=LectureNavy}
\setbeamercolor{palette primary}{bg=LectureNavy,fg=white}
\setbeamercolor{palette secondary}{bg=LectureAccent,fg=white}
\setbeamercolor{palette tertiary}{bg=LectureGray,fg=white}
% Madrid title box is dark --- title/subtitle must be light for contrast
\setbeamercolor{title}{fg=white,bg=LectureNavy}
\setbeamercolor{subtitle}{fg=white,bg=LectureNavy}
\setbeamercolor{author}{fg=LectureNavy}
\setbeamercolor{date}{fg=LectureGray}
\setbeamercolor{institute}{fg=LectureGray}
\setbeamercolor{frametitle}{bg=LectureNavy,fg=white}
\setbeamercolor{block title}{bg=LectureNavy,fg=white}
\setbeamercolor{block body}{bg=CodeBg,fg=black}
\setbeamercolor{block title alerted}{bg=LectureAccent,fg=white}
\setbeamercolor{block body alerted}{bg=TryBg,fg=black}
\setbeamercolor{block title example}{bg=LectureGray,fg=white}
\setbeamercolor{block body example}{bg=AlertBg,fg=black}

\setbeamertemplate{navigation symbols}{}
\setbeamertemplate{itemize items}[circle]
\setbeamertemplate{enumerate items}[default]

\usepackage[T1]{fontenc}
\usepackage[utf8]{inputenc}
\usepackage{lmodern}
\usepackage{booktabs}
\usepackage{tabularx}
\usepackage{graphicx}
\usepackage{hyperref}
\usepackage{listings}
\usepackage{xcolor}
\usepackage{tikz}
\usetikzlibrary{positioning,arrows.meta,shapes.geometric}

% listings: Python without shell-escape
\lstdefinestyle{pythonstyle}{
    language=Python,
    basicstyle=\ttfamily\small,
    keywordstyle=\color{KeywordBlue}\bfseries,
    commentstyle=\color{CommentGray}\itshape,
    stringstyle=\color{StringGreen},
    numbers=left,
    numberstyle=\tiny\color{CommentGray},
    numbersep=6pt,
    frame=single,
    rulecolor=\color{LectureGray!40},
    backgroundcolor=\color{CodeBg},
    breaklines=true,
    showstringspaces=false,
    tabsize=4,
    captionpos=b,
}
\lstset{style=pythonstyle}

\newcommand{\pycode}[1]{{\ttfamily\small\detokenize{#1}}}
\newcommand{\trythis}[1]{%
    \begin{alertblock}{Try this}
    #1
    \end{alertblock}%
}
\newcommand{\commonmistake}[1]{%
    \begin{exampleblock}{Common mistake}
    #1
    \end{exampleblock}%
}
\newcommand{\remember}[1]{%
    \begin{block}{Remember}
    #1
    \end{block}%
}

% Formative in-class checks: question slide, then a dedicated Answer slide
\newcommand{\quizpause}{%
    \par\vspace{0.5em}
    {\small\textit{Pause --- answer (clicker / raise hand / neighbour) before the next slide.}}%
}
\newcommand{\quizbox}[1]{%
    \setbeamercolor{block title}{bg=KeywordBlue,fg=white}%
    \setbeamercolor{block body}{bg=QuizBg,fg=black}%
    \begin{block}{Quiz}
    #1
    \end{block}%
    \setbeamercolor{block title}{bg=LectureNavy,fg=white}%
    \setbeamercolor{block body}{bg=CodeBg,fg=black}%
}
\newcommand{\answerbox}[1]{%
    \setbeamercolor{block title}{bg=AnswerGold,fg=white}%
    \setbeamercolor{block body}{bg=AnswerKeyBg,fg=black}%
    \begin{block}{Answer}
    #1
    \end{block}%
    \setbeamercolor{block title}{bg=LectureNavy,fg=white}%
    \setbeamercolor{block body}{bg=CodeBg,fg=black}%
}

% Empty author/date placeholders for the lecturer to fill
\author{}
\date{}


\usetheme{Madrid}
\usecolortheme{default}

% Clean palette: deep slate + teal accent (distinct from Java navy decks)
\definecolor{LectureNavy}{HTML}{1A365D}
\definecolor{LectureAccent}{HTML}{319795}
\definecolor{LectureGray}{HTML}{4A5568}
\definecolor{CodeBg}{HTML}{F7FAFC}
\definecolor{AlertBg}{HTML}{FFF5F5}
\definecolor{TryBg}{HTML}{F0FFF4}
\definecolor{QuizBg}{HTML}{EBF8FF}
\definecolor{AnswerKeyBg}{HTML}{FFFFF0}
\definecolor{AnswerGold}{HTML}{B7791F}
\definecolor{KeywordBlue}{HTML}{2B6CB0}
\definecolor{StringGreen}{HTML}{276749}
\definecolor{CommentGray}{HTML}{718096}

\setbeamercolor{structure}{fg=LectureNavy}
\setbeamercolor{palette primary}{bg=LectureNavy,fg=white}
\setbeamercolor{palette secondary}{bg=LectureAccent,fg=white}
\setbeamercolor{palette tertiary}{bg=LectureGray,fg=white}
% Madrid title box is dark --- title/subtitle must be light for contrast
\setbeamercolor{title}{fg=white,bg=LectureNavy}
\setbeamercolor{subtitle}{fg=white,bg=LectureNavy}
\setbeamercolor{author}{fg=LectureNavy}
\setbeamercolor{date}{fg=LectureGray}
\setbeamercolor{institute}{fg=LectureGray}
\setbeamercolor{frametitle}{bg=LectureNavy,fg=white}
\setbeamercolor{block title}{bg=LectureNavy,fg=white}
\setbeamercolor{block body}{bg=CodeBg,fg=black}
\setbeamercolor{block title alerted}{bg=LectureAccent,fg=white}
\setbeamercolor{block body alerted}{bg=TryBg,fg=black}
\setbeamercolor{block title example}{bg=LectureGray,fg=white}
\setbeamercolor{block body example}{bg=AlertBg,fg=black}

\setbeamertemplate{navigation symbols}{}
\setbeamertemplate{itemize items}[circle]
\setbeamertemplate{enumerate items}[default]

\usepackage[T1]{fontenc}
\usepackage[utf8]{inputenc}
\usepackage{lmodern}
\usepackage{booktabs}
\usepackage{tabularx}
\usepackage{graphicx}
\usepackage{hyperref}
\usepackage{listings}
\usepackage{xcolor}
\usepackage{tikz}
\usetikzlibrary{positioning,arrows.meta,shapes.geometric}

% listings: Python without shell-escape
\lstdefinestyle{pythonstyle}{
    language=Python,
    basicstyle=\ttfamily\small,
    keywordstyle=\color{KeywordBlue}\bfseries,
    commentstyle=\color{CommentGray}\itshape,
    stringstyle=\color{StringGreen},
    numbers=left,
    numberstyle=\tiny\color{CommentGray},
    numbersep=6pt,
    frame=single,
    rulecolor=\color{LectureGray!40},
    backgroundcolor=\color{CodeBg},
    breaklines=true,
    showstringspaces=false,
    tabsize=4,
    captionpos=b,
}
\lstset{style=pythonstyle}

\newcommand{\pycode}[1]{{\ttfamily\small\detokenize{#1}}}
\newcommand{\trythis}[1]{%
    \begin{alertblock}{Try this}
    #1
    \end{alertblock}%
}
\newcommand{\commonmistake}[1]{%
    \begin{exampleblock}{Common mistake}
    #1
    \end{exampleblock}%
}
\newcommand{\remember}[1]{%
    \begin{block}{Remember}
    #1
    \end{block}%
}

% Formative in-class checks: question slide, then a dedicated Answer slide
\newcommand{\quizpause}{%
    \par\vspace{0.5em}
    {\small\textit{Pause --- answer (clicker / raise hand / neighbour) before the next slide.}}%
}
\newcommand{\quizbox}[1]{%
    \setbeamercolor{block title}{bg=KeywordBlue,fg=white}%
    \setbeamercolor{block body}{bg=QuizBg,fg=black}%
    \begin{block}{Quiz}
    #1
    \end{block}%
    \setbeamercolor{block title}{bg=LectureNavy,fg=white}%
    \setbeamercolor{block body}{bg=CodeBg,fg=black}%
}
\newcommand{\answerbox}[1]{%
    \setbeamercolor{block title}{bg=AnswerGold,fg=white}%
    \setbeamercolor{block body}{bg=AnswerKeyBg,fg=black}%
    \begin{block}{Answer}
    #1
    \end{block}%
    \setbeamercolor{block title}{bg=LectureNavy,fg=white}%
    \setbeamercolor{block body}{bg=CodeBg,fg=black}%
}

% Empty author/date placeholders for the lecturer to fill
\author{}
\date{}


\usetheme{Madrid}
\usecolortheme{default}

% Clean palette: deep slate + teal accent (distinct from Java navy decks)
\definecolor{LectureNavy}{HTML}{1A365D}
\definecolor{LectureAccent}{HTML}{319795}
\definecolor{LectureGray}{HTML}{4A5568}
\definecolor{CodeBg}{HTML}{F7FAFC}
\definecolor{AlertBg}{HTML}{FFF5F5}
\definecolor{TryBg}{HTML}{F0FFF4}
\definecolor{QuizBg}{HTML}{EBF8FF}
\definecolor{AnswerKeyBg}{HTML}{FFFFF0}
\definecolor{AnswerGold}{HTML}{B7791F}
\definecolor{KeywordBlue}{HTML}{2B6CB0}
\definecolor{StringGreen}{HTML}{276749}
\definecolor{CommentGray}{HTML}{718096}

\setbeamercolor{structure}{fg=LectureNavy}
\setbeamercolor{palette primary}{bg=LectureNavy,fg=white}
\setbeamercolor{palette secondary}{bg=LectureAccent,fg=white}
\setbeamercolor{palette tertiary}{bg=LectureGray,fg=white}
% Madrid title box is dark --- title/subtitle must be light for contrast
\setbeamercolor{title}{fg=white,bg=LectureNavy}
\setbeamercolor{subtitle}{fg=white,bg=LectureNavy}
\setbeamercolor{author}{fg=LectureNavy}
\setbeamercolor{date}{fg=LectureGray}
\setbeamercolor{institute}{fg=LectureGray}
\setbeamercolor{frametitle}{bg=LectureNavy,fg=white}
\setbeamercolor{block title}{bg=LectureNavy,fg=white}
\setbeamercolor{block body}{bg=CodeBg,fg=black}
\setbeamercolor{block title alerted}{bg=LectureAccent,fg=white}
\setbeamercolor{block body alerted}{bg=TryBg,fg=black}
\setbeamercolor{block title example}{bg=LectureGray,fg=white}
\setbeamercolor{block body example}{bg=AlertBg,fg=black}

\setbeamertemplate{navigation symbols}{}
\setbeamertemplate{itemize items}[circle]
\setbeamertemplate{enumerate items}[default]

\usepackage[T1]{fontenc}
\usepackage[utf8]{inputenc}
\usepackage{lmodern}
\usepackage{booktabs}
\usepackage{tabularx}
\usepackage{graphicx}
\usepackage{hyperref}
\usepackage{listings}
\usepackage{xcolor}
\usepackage{tikz}
\usetikzlibrary{positioning,arrows.meta,shapes.geometric}

% listings: Python without shell-escape
\lstdefinestyle{pythonstyle}{
    language=Python,
    basicstyle=\ttfamily\small,
    keywordstyle=\color{KeywordBlue}\bfseries,
    commentstyle=\color{CommentGray}\itshape,
    stringstyle=\color{StringGreen},
    numbers=left,
    numberstyle=\tiny\color{CommentGray},
    numbersep=6pt,
    frame=single,
    rulecolor=\color{LectureGray!40},
    backgroundcolor=\color{CodeBg},
    breaklines=true,
    showstringspaces=false,
    tabsize=4,
    captionpos=b,
}
\lstset{style=pythonstyle}

\newcommand{\pycode}[1]{{\ttfamily\small\detokenize{#1}}}
\newcommand{\trythis}[1]{%
    \begin{alertblock}{Try this}
    #1
    \end{alertblock}%
}
\newcommand{\commonmistake}[1]{%
    \begin{exampleblock}{Common mistake}
    #1
    \end{exampleblock}%
}
\newcommand{\remember}[1]{%
    \begin{block}{Remember}
    #1
    \end{block}%
}

% Formative in-class checks: question slide, then a dedicated Answer slide
\newcommand{\quizpause}{%
    \par\vspace{0.5em}
    {\small\textit{Pause --- answer (clicker / raise hand / neighbour) before the next slide.}}%
}
\newcommand{\quizbox}[1]{%
    \setbeamercolor{block title}{bg=KeywordBlue,fg=white}%
    \setbeamercolor{block body}{bg=QuizBg,fg=black}%
    \begin{block}{Quiz}
    #1
    \end{block}%
    \setbeamercolor{block title}{bg=LectureNavy,fg=white}%
    \setbeamercolor{block body}{bg=CodeBg,fg=black}%
}
\newcommand{\answerbox}[1]{%
    \setbeamercolor{block title}{bg=AnswerGold,fg=white}%
    \setbeamercolor{block body}{bg=AnswerKeyBg,fg=black}%
    \begin{block}{Answer}
    #1
    \end{block}%
    \setbeamercolor{block title}{bg=LectureNavy,fg=white}%
    \setbeamercolor{block body}{bg=CodeBg,fg=black}%
}

% Empty author/date placeholders for the lecturer to fill
\author{}
\date{}
